%% kapa: exam=1 round=35 year=2024 period=2 form=A subject=accounting subject_ko=회계학 seq=34 number=74
%% source: Q-Net 감정평가사 기출문제 (cst00309) artlSeq=5213229 file=raw/1008/5213229/2교시 문제지 원본.pdf
%% conversion: kordoc 4.24.51 -> tools/gen1.py (auto); review=proofread-sonnet; flags=none
\begin{question}{74}
(주)감평은 20×1년 초 영업을 개시하였으며, 표준원가계산제도를 채택하고 있다. 직접재료 kg당 실제 구입가격은 ￦5, 제품 단위당 직접재료 표준원가는 ￦6(2kg×￦3/kg)이다. 직접재료원가에 대한 차이 분석결과 구입가격차이가 ￦3,000(불리), 능률차이가 ￦900(유리)이다. 20×1년 실제 제품 생산량이 800단위일 때, 기말 직접재료 재고수량은? (단, 기말재공품은 없다.)
\begin{choices}
  \item 50kg
  \item 100kg
  \item 130kg
  \item 200kg
  \item 230kg
\end{choices}
\end{question}
\answer{4}
